\chapter{Post-Training Generalisation}
\label{chap:Downstream Post-Training Generalisation}
Often, a model is trained (at great cost) and needs to either \textit{specialise} to a subsection of the data, which was under-represented in the training, or \textit{generalise} to a new, unseen set of data.
When designing this kind of procedure, we need to ask ourselves \textit{what is shared, what changed, and how much of the model should be allowed to adapt?}


\section{Transfer Learning}
\label{sec:adaptation-transfer-learning-reuse-what-the-model-has-learned}

When \textit{transferring} to a new task, the main design choice is how much of the pretrained model we allow to change.

\begin{calloutmemoryhook}
    In post-training generalisation, overfitting is basically forgetting.
\end{calloutmemoryhook}

A model is \textbf{tuned} when training is continued from where the main training ended, normally with some carefully curated dataset.
For example, an image model may be found online, and further trained on a curated dataset of images in the style of Disney Animations, leading to a model which has \textit{specialised} in Disney's style, while retaining a lot of the general image model capabilities of the parent model.

A \textbf{Prediction Head} is a small layer on the end of a model which is added to handle a specific task. 
For example, imagine training a model to predict house prices and then wanting to predict the probability of a house selling in the next year. 
The problems share a lot of information, and so it may be worth just replacing the final layer of the house-price model and training just that layer on the new probability-of-sale data.

\textbf{Negative transfer} means that your post-trained model performs worse than your pre-trained model.

\begin{calloutmemoryhook}
    The word \textit{Fine-Tuning} should never have been invented, it just means \textit{continued training on custom data}.
\end{calloutmemoryhook}

\begin{table}[htbp]
    \centering
    \small
    \begin{tabularx}{\linewidth}{@{}L{.19\linewidth}YYY@{}}
    \toprule
    Strategy & What changes? & Main trade-off & Applications \\
    \midrule

    \textbf{Frozen}
    &
    The main model is kept the same, and a small secondary model is duct-taped onto the end, creating a specialised \textit{prediction head}.
    &
    Very cheap to train and retains the original representations, but the main model cannot adapt to the new task.
    &
    Small datasets where the pretrained features are already useful. E.g.\ adding a classifier to a pretrained image encoder.
    \\

    \addlinespace
    \textbf{Fully Tuned}
    &
    The entire model's training is continued on the new dataset. A prediction head is optional.
    &
    Maximum flexibility but highest cost, greater risk of overfitting.
    &
    Large, specialised datasets or substantial domain shifts. E.g.\ adapting a general image model to medical imaging.
    \\

    \addlinespace
    \textbf{Partially Tuned}
    &
    A prediction head is added, and a subset of the pretrained layers are also further trained, allowing the main model to specialise.
    &
    Balances adaptation against training cost, but performance depends on which layers are unfrozen.
    &
    Tasks where general low-level features remain useful but higher-level representations need adapting. E.g.\ retraining the final CNN blocks for a new classification task.
    \\

    \addlinespace
    \textbf{Efficiently Tuned}
    &
    Special, highly efficient methods of continuing training which only train a small number of parameters.
    &
    Much lower trainable-parameter and optimiser-memory cost, but constrains how the model can adapt. Total compute savings depend on the method.
    &
    Adapting large models under limited memory or maintaining many specialised versions of one base model. E.g.\ LoRA fine-tuning of an LLM.
    \\

    \bottomrule
    \end{tabularx}
    \caption{Transfer learning is a continuum of adaptation capacity, trading training cost against the freedom to specialise pretrained representations.}
\end{table}


\subsection{Domain Adaptation}
\label{sec:adaptation-domain-adaptation-the-data-distribution-changes}
It is useful to think of domain adatation as \textit{which probability distribution (or combination) changed?}
\begin{align}
    \underbrace{P_S(x)\ne P_T(x),\;P_S(y\mid x)=P_T(y\mid x)}_{\text{covariate shift: inputs change, rule stays}},\\
    \underbrace{P_S(y)\ne P_T(y),\;P_S(x\mid y)=P_T(x\mid y)}_{\text{label shift: class prevalence changes}},\\
    \underbrace{P_S(y\mid x)\ne P_T(y\mid x)}_{\text{concept shift: prediction rule changes}}.
\end{align}

Under covariate shift, examples can be \textbf{reweighted} to represent the target input distribution.
The input data itself can be adapted to fit the new distribution, like
\[
    \underbrace{\E_{P_T}[\ell]}_{\text{target risk}}
    =
    \underbrace{\E_{P_S}\!\left[
    \frac{P_T(x)}{P_S(x)}\ell(f(x),y)
    \right]}_{\text{source risk weighted toward target-like inputs}}.
\]
However, for small target regions (small classes, sparse prediction space, etc.) training on reweighted data just causes overfitting.

\figpdf[.8\textwidth]{domain_shift}{Synthetic covariate shift in one input coordinate. This feature could be reweighted as the distribution is just a transform (changed from Normal to another Normal).}

\begin{calloutmemoryhook}
    Reweighting can fix input-frequency shift but it cannot fix a changed input-to-label rule, and breaks under class imbalance.
\end{calloutmemoryhook}

\subsection{Summary}
\begin{calloutbias}
    \begin{tabularx}{\linewidth}{@{}L{.23\linewidth}YY@{}}
        \toprule
        Approach & Shared structure assumed & Failure to test \\
        \midrule
        Frozen transfer & Source representation retains target signal & Target feature was suppressed upstream \\
        Fine-tuning & Pretrained parameters are a useful starting point & Small target set overfits \\
        Domain invariance & Some task-relevant features persist across domains & Alignment mixes different classes \\
        Prototypical networks & Classes form meaningful embedding centroids & Classes are multimodal/overlapping \\
        MAML & New tasks resemble meta-training tasks & Task family changes substantially \\
        Zero-shot multimodal & Modalities share semantic geometry & Prompt/class semantics are ambiguous \\
        LoRA & Adaptation is approximately low rank & Task needs broad independent parameter changes \\
        \bottomrule
    \end{tabularx}
\end{calloutbias}

\section{Specialise: Zero-shot learning}
\label{sec:adaptation-zero-shot-learning-no-task-specific-labelled-examples}
Zero-shot learning is only really available to Large Language/Image Models, because few other models generalise as well as them.
It is done by changing the context that the model has available, rather than by changing weights.
For a shared image/text embedding, classification can be ``choose the text description with highest cosine similarity'':
\[
    \hat y
    =
    \argmax_j
    \underbrace{
    \frac{f_\theta(x)^\top g_\phi(t_j)}{\|f_\theta(x)\|_2\|g_\phi(t_j)\|_2}
    }_{\text{image--text directional similarity}}.
\]
Performance depends on the pre-training semantics, candidate class set and prompt wording.
Zero-shot does not mean zero prior knowledge, it assumes that the specific application exists within the pre-training dataset, and all we need to do is \textit{prompt better}.

\begin{calloutmemoryhook}
    Prompting changes context, while fine-tuning changes persistent parameters.
\end{calloutmemoryhook}


\section{Specialise: Few-shot learning}
\label{sec:adaptation-few-shot-learning-learn-from-very-few-examples}
Few shot learning allows us to curate a small dataset and conduct further training with it.
In the usual ``N-way K-shot'' terminology, an episode contains a small support set for each candidate class plus query examples for evaluation.
Few-shot learning is an approach, not a learning mechanism, so it can mean metric learning, gradient-based adaptation, frozen embeddings, or in-context prompting.

We reduce each class to a centroid:
\[
    \underbrace{c_{\text{class }j}}_{\text{class prototype}}
    =
    \frac{1}{N_{\text{support per class}}}
    \sum_{x_i\in S_j}f_\theta(x_i),
    \qquad
    \hat y_{\text{query}}
    =
    \argmin_j d\!\left(f_\theta(x_{\text{query}}),c_{\text{class }j}\right).
\]
This only works if the learned geometry makes same-class examples cluster in a transferable way.
\figpdf[.8\textwidth]{few_shot_prototypes}{Few labelled support examples define class prototypes; a query is classified by embedding-space proximity.}

\section{Generalise: Meta-learning and MAML}
\label{sec:adaptation-meta-learning-and-maml}
Meta-learning trains across tasks so that a new related task can be adapted quickly.
MAML makes that idea explicit by optimising query performance \emph{after} a task-specific support-set update:
\[
    \min_\theta
    \sum_{j=1}^{N_{\text{tasks}}}
    \mathcal L^{\text{query}}_j\!\left(
    \underbrace{
    \theta-\alpha\nabla_\theta\mathcal L^{\text{support}}_j(\theta)
    }_{\theta'_j:\ \text{one task-adaptation step}}
    \right).
\]
Exact differentiation through the inner update introduces second-order dependence.
First-order variants drop that cost.
The Hessian/second-order connection is in \cref{subsec:la-hessian,sec:optimisation-newton-s-method-and-other-second-order-methods}.



\section{Specialise: Parameter Efficient Continued-Training}
The general idea with Parameter Efficient Fine-Tuning (PEFT) is to avoid changing the large model instance.
Activations, optimiser states and the other backpropagation objects are very large to store in memory, and most specialisation problems do not require scale.
It's wasteful to train a full set of 1 billion parameters on 2000 examples!
So, we transform the weights into a low dimension using various transforming methods. 
We work in that lower dimensional space and then un-transform the results into the model.

\subsection{Low Rank Adapters (LoRA)}
\label{sec:adaptation-lora-a-parameter-efficient-adaptation-mechanism}
LoRA freezes a pretrained weight matrix and learns a low-rank update:
\[
    \underbrace{\theta'_{\text{weights}}}_{\text{adapted matrix}}
    =
    \underbrace{\theta_{\text{weights}}}_{\text{frozen pretrained matrix}}
    +
    \underbrace{\theta_{\text{up}}\theta_{\text{down}}}_{\Delta\theta:\ \text{low-rank update}}.
\]
This changes the bottleneck to the cost of dimension reduction/expansion:
\[
    \underbrace{\theta_{\text{down}}}_{\text{compress}}
    \in\R^{N_{\text{adaptation directions}}\times N_{\text{input features}}},
    \qquad
    \underbrace{\theta_{\text{up}}}_{\text{expand}}
    \in\R^{N_{\text{output features}}\times N_{\text{adaptation directions}}}.
\]
The parameter saving is:
\[
    \underbrace{N_{\text{adaptation directions}}
    \left(N_{\text{input features}}+N_{\text{output features}}\right)}_{\text{trainable LoRA parameters}}
    \quad\text{vs.}\quad
    \underbrace{N_{\text{input features}}N_{\text{output features}}}_{\text{full matrix parameters}}.
\]
The underlying assumption is that useful task adaptation lives in relatively few independent directions.
Low trainable-parameter count does not imply every training/inference cost shrinks proportionally.


\subsection{Quantised Low Rank Adapters (QLoRA)}
\label{subsec:adaptation-qlora}
QLoRA combines two different restrictions: keep the pretrained base model in a low-bit representation to save memory, while training LoRA adapters for task-specific change.
Conceptually,
\[
    \underbrace{\text{quantised frozen base}}_{\text{cheap storage}}
    +
    \underbrace{\text{trainable low-rank update}}_{\text{adaptation capacity}}.
\]
Quantisation and low-rank adaptation solve different bottlenecks, so they compose naturally.
Quantisation mechanics and their quality--memory trade-off are developed in \cref{sec:inference-quantisation}.

\begin{calloutmemoryhook}
    QLoRA compresses the frozen base, and LoRA restricts the trainable change.
\end{calloutmemoryhook}


